% NeurIPS 2026 Workshop Submission (4-page anonymous)
% "Prediction Bottlenecks Don't Discover Causal Structure"
% Trimmed from main.tex for the 4-page anonymous format.

\documentclass{article}
\usepackage{../neurips_2026}
\usepackage[utf8]{inputenc}
\usepackage[T1]{fontenc}
\usepackage{hyperref}
\usepackage{url}
\usepackage{booktabs}
\usepackage{amsfonts}
\usepackage{nicefrac}
\usepackage{microtype}
\usepackage{graphicx}
\usepackage{xcolor}
\usepackage{amsmath}
\usepackage{amssymb}
\usepackage{enumitem}

\title{Prediction Bottlenecks Don't Discover Causal Structure (But Here's What They Actually Do)}

\author{%
  Anonymous Author(s)\\
  Affiliation\\
  Address\\
  \texttt{email}\\
}

\begin{document}
\maketitle

\begin{abstract}
A Mamba state-space model trained only for next-step prediction appears to
recover Granger-causal structure through a strikingly simple readout:
$\mathbf{S} = |\mathbf{W}_{\text{out}} \mathbf{W}_{\text{in}}|$. Early
experiments suggested this generalized across architectures and benefited
from interventional data with $p < 10^{-5}$. We walk this claim through a
five-stage falsification pipeline. The method-level claim does not survive:
(i) a plain linear bottleneck does as well or better, (ii) tuned Lasso
beats it on synthetic and real data (NOAA climate, ETFs, Lorenz-96), and
(iii) the headline intervention advantage is $\sim$60\% a size-matched-data
confound and $\sim$40\% a target-corruption artifact specific to a
non-standard random-forcing scheme. What survives is a narrow
characterization result. We hope the falsification scaffold is as useful
to the community as the positive findings we could not salvage.
\end{abstract}

% ============================================================
\section{Introduction}
% ============================================================

A natural hope in time-series modelling is that a forecaster trained only
on next-step prediction will, as a byproduct, expose causal or dependency
structure through its learned weights — \emph{causal discovery as a
byproduct of prediction}. If true, every pretrained forecaster would
double as a causal-discovery engine.

This paper reports the result of taking that hope seriously. Our starting
point was the observation that for a Mamba selective state-space
model~\citep{gu2023mamba} trained for next-step MSE, the product
$\mathbf{S} = |\mathbf{W}_{\text{out}} \mathbf{W}_{\text{in}}|$ recovers
Granger-causal graphs with striking reliability on synthetic VAR
benchmarks. Initial experiments suggested the phenomenon generalized
across architectures, emerged from training rather than initialization,
and — most provocatively — that interventional data amplified the
bottleneck's advantage with $p < 10^{-5}$ across twelve $(K, T)$
configurations.

Each of these claims turned out to be less robust than it first appeared.
In the spirit of ICBINB, our contribution is \emph{negative}: we show
that a natural method candidate fails in five subtly different ways,
identify the controls that caused the failures, and report the narrow
findings that remain valid.

\paragraph{Contributions.}
\begin{enumerate}[leftmargin=1.5em, topsep=0pt, itemsep=0pt]
  \item We falsify architecture specificity: a linear bottleneck matches
  or beats the SSM under matched capacity and seeds.
  \item We falsify competitiveness with sparse regression: tuned Lasso
  beats the bottleneck on synthetic CauseMe-style benchmarks and on real
  observational data (NOAA climate, sector ETFs, Lorenz-96).
  \item We falsify the headline intervention advantage: under a
  size-matched control and proper $\text{do}(X = c)$ interventions, the
  effect collapses; the residual reproduces (and is larger) under
  classical Granger.
  \item We document what survives — mild-nonlinearity, sample-efficiency,
  target-corruption robustness — as characterization, not method, claims,
  and release a reusable falsification scaffold.
\end{enumerate}

% ============================================================
\section{Method: The Extraction We Are Falsifying}
\label{sec:observation}
% ============================================================

For a $K$-variable series $\mathbf{X} \in \mathbb{R}^{T \times K}$, we
train a model with input projection $\mathbf{W}_{\text{in}} \in
\mathbb{R}^{d \times K}$ and output projection $\mathbf{W}_{\text{out}}
\in \mathbb{R}^{K \times d}$ on next-step MSE. No causal-specific loss.
After training, we extract
$\mathbf{S} = |\mathbf{W}_{\text{out}} \mathbf{W}_{\text{in}}| \in
\mathbb{R}^{K \times K}$ with $\mathbf{S}_{i,j}$ interpreted as the
strength of $j \to i$, zero the diagonal, normalize, and threshold. For
architectures with explicit lagged input (the lagged bottleneck variant),
we extend to $\mathbf{S}^{(\tau)} = |\mathbf{W}_{\text{out}}
\mathbf{W}^{(\tau)}_{\text{in}}|$ and evaluate via flat-lag AUROC against
ground-truth adjacency.

% ============================================================
\section{Five Falsifications}
\label{sec:falsifications}
% ============================================================

\paragraph{F1. Architecture does not matter.} A pre-registered sweep at
matched parameter counts and seeds across four datasets and five
architectures (Mamba SSM, Transformer, LSTM, MLP, Linear) finds the
linear bottleneck matches or beats the SSM on every dataset (mean
AUROC: VAR(1)-chain $K{=}5$: SSM $1.00$, Linear $1.00$;
VAR(1)-random $K{=}10$: SSM $0.77$, Linear $\mathbf{0.99}$;
regime-switch $K{=}3$: both $1.00$;
Lorenz $K{=}3$: SSM $0.52$, Linear $\mathbf{0.94}$). The
``emergence'' is the linear bottleneck doing low-rank lag-1 regression;
the SSM achieves the same estimator at higher parameter cost.

\paragraph{F2. Tuned Lasso beats the bottleneck on graph recovery.}
Across a focused stress grid ($K \in \{10, 20\}$, $T \in \{150, 300\}$,
max-lag 8, Table~\ref{tab:lasso}) the bottleneck wins only $12\%$ of
graph-recovery runs and $0\%$ of MSE runs against tuned baselines. On
tigramite CauseMe-style processes the gap widens with $K$: Lasso reaches
AUROC $0.98$ at $K{=}20$ versus $0.73$ for the bottleneck.

\begin{table}[t]
\centering
\small
\caption{Bottleneck vs.\ best tuned baseline on the focused stress grid.}
\label{tab:lasso}
\begin{tabular}{lrr}
\toprule
Metric vs.\ \{OLS, Ridge, Lasso, RRR\} & Mean AUROC delta & Win rate \\
\midrule
Bottleneck flat-lag AUROC      & $-0.050$ & $0.12$ \\
Bottleneck prediction MSE      & $+0.286$ & $0.00$ \\
Best baseline by AUROC: Lasso ($38/48$), RRR ($10/48$) & & \\
Best baseline by MSE:   Lasso ($48/48$)              & & \\
\bottomrule
\end{tabular}
\end{table}

\paragraph{F3. Real data: not competitive.} We evaluate six methods on
three datasets of progressively cleaner ground truth (Table~\ref{tab:real}):
63 years of NOAA monthly climate indices (ENSO, NAO, PDO, AMO, SOI, PNA)
with teleconnection edges from the climate
literature~\citep{trenberth1994decadal,delworth2000observed,newman2016pdo};
7.5 years of daily SPDR sector ETF log returns ($K{=}10$, soft labels);
and Lorenz-96 ($K{=}10$, $F{=}10$) with mathematically unambiguous
ground truth~\citep{tank2021neural,nauta2019causal}. We deliberately
exclude the ENSO$\leftrightarrow$SOI pair from climate ground truth
because SOI and the Niño 3.4 SST anomaly are the same physical signal
with opposite sign, not a causal edge. On Lorenz-96, where ground truth
is unambiguous, \emph{PCMCI and bivariate Granger dominate} (0.986,
0.979); the lagged bottleneck trails at $0.920$ and the SSM bottleneck
is worst at $0.802$. The bottleneck is not the best method on any
dataset with clean ground truth.

\begin{table}[t]
\centering
\small
\caption{AUROC on three real/benchmark datasets. The bottleneck tops only
climate, where ground truth is weak (3 edges, $K{=}6$); on the clean
Lorenz-96 benchmark, classical causal methods lead by a wide margin.}
\label{tab:real}
\begin{tabular}{lccc}
\toprule
Method & Climate (NOAA) & Finance (ETFs) & Lorenz-96 \\
 & $K{=}6, T{=}757$ & $K{=}10, T{=}1893$ & $K{=}10, T{=}1500$ \\
\midrule
Weight-Proj (SSM) & $\mathbf{0.864}$ & $0.438$ & $0.802$ \\
PCMCI & $0.704$ & $0.510$ & $\mathbf{0.986}$ \\
Tuned Lasso & $0.630$ & $0.500$ & $0.974$ \\
Lagged Bottleneck & $0.617$ & $0.514$ & $0.920$ \\
Tuned Ridge & $0.617$ & $\mathbf{0.696}$ & $0.974$ \\
Bivariate Granger & $0.605$ & $0.498$ & $0.979$ \\
\bottomrule
\end{tabular}
\end{table}

\paragraph{F4. Interventions: a confound story.} An earlier experiment
(Exp 19) compared $\text{AUROC}(\mathbf{X}_{\text{obs}})$ versus
$\text{AUROC}(\mathbf{X}_{\text{obs}} \cup \mathbf{X}_{\text{int}})$ for
both bottleneck and Lasso, reporting $12/12$ configs significant at
$p < 10^{-5}$ in favor of the bottleneck. Two controls dissolve it.
\emph{Size-match (Exp 22):} Exp 19 conflated more data with intervention
content; adding a fresh observational arm of equal total size collapses
the effect to $\sim$40\% of the raw gap. \emph{Intervention type
(Exp 23):} the original ``intervention'' was a per-step random forcing
$x_{i,t} \leftarrow s\cdot \epsilon_{i,t}$, not a standard
$\text{do}$-intervention. Replacing it with constant clamps
$\text{do}(X_i = c)$ gives only $3/12$ configs significant (mean gap
$+0.002$); soft (added-noise) interventions give $6/12$, mean gap
$+0.008$. Only random forcing produces a stable bottleneck gain, best
explained as robustness to corrupted target rows.

\paragraph{F5. Method-agnosticity.} The residual effect is not unique to
the bottleneck. Adding PCMCI and bivariate Granger to the size-matched
control: the bottleneck shows a significant size-matched gain ($+0.026$
to $+0.054$, $p < 10^{-3}$), \emph{but so does Granger}, with larger
effects at higher $K$ ($+0.040$ at $K{=}20$, $+0.095$ at $K{=}30$).
Lasso uniquely fails, consistent with per-equation regression breaking
under target corruption. The intervention effect, insofar as it survives
controls, is method-agnostic and matches a target-corruption-robustness
story rather than causal-structure extraction.

% ============================================================
\section{What Survives}
\label{sec:survives}
% ============================================================

After five falsifications, three narrow positives remain.
\textbf{(i) Mild-nonlinearity regime:} in a regime sweep over noise,
density, and nonlinearity strength, the bottleneck outperforms Lasso at
\texttt{nonlinear=0.3} ($78\%$ win rate over 3 seeds and 9 hyperparameter
cells). At stronger nonlinearity, both methods fail equally.
\textbf{(ii) Sample efficiency:} bottleneck gains from additional
observational data exceed Lasso's by $\sim 0.07$ AUROC at $K \in \{20,
30\}$ under size-matched control arms — modest but useful where collecting
more observations is cheap.
\textbf{(iii) Target-corruption robustness:} the bottleneck's shared
$\mathbf{W}_{\text{out}}$ absorbs per-step random forcing better than
Lasso's per-equation fits. This explains the residual intervention
effect, but is a reliability result, not a causal-discovery one.

These are characterization claims. We do not propose the bottleneck as a
causal discovery method.

% ============================================================
\section{Lessons}
% ============================================================

\paragraph{Size-matched controls beat seed counts.} Exp 19 had 5 seeds.
Exp 21 had 20 seeds and reported $p < 10^{-5}$ across 12 configs. Both
found the same effect; neither would have flagged it as a sample-size
confound. Adding seeds tightens a noisy measurement; it cannot diagnose
a missing structural control.

\paragraph{Pre-register the intervention scheme.} The gap between
``random forcing'' and proper $\text{do}(X_i = c)$ is the difference
between a significant effect and a null, and it is not obvious from the
code. We should have written down which intervention semantics counted
as ``interventional data'' before running the experiment.

\paragraph{Prediction baselines are not causal baselines.} Our
intervention experiments compared the bottleneck only against
prediction-fit baselines (OLS, Ridge, Lasso, RRR). It was natural to
miss that classical causal methods (Granger, PCMCI) would benefit from
the same intervention scheme by the same mechanism. Any future
``neural causal discovery'' claim should include at least one classical
causal baseline in the intervention control arm.

% ============================================================
\section{Conclusion}
% ============================================================

We set out to find a free causal-discovery method hiding inside
prediction bottlenecks and could not. The headline effects either fail
under size-matched controls, are specific to a non-standard intervention
scheme, or reproduce as larger effects in classical causal methods. What
survives is a narrow set of characterization findings and a reusable
falsification scaffold.

\paragraph{Code.} All experiments, data loaders, and a one-line
reproduction script will be released upon acceptance (anonymous URL
provided to reviewers).

\bibliographystyle{plainnat}
\bibliography{../references}

\end{document}
